\hypertarget{appendix-a1b1b4580dfd}{}
\section{Floating-point addition and absorption}\label{appendix-a1b1b4580dfd-floating-point-addition-and-absorption}

\textbf{Curriculum depth:} Introduced · \textbf{First seen:} C005 · \textbf{Latest development:} C005

\subsection{Definition and mechanism}\label{appendix-a1b1b4580dfd-definition-and-mechanism}

Addition aligns significands to a common exponent before adding/subtracting and rounding. At 2\^{}24, upward binary32 spacing is 2, so adding 1 is a tie absorbed to the same even endpoint. 1.5+0.15625 aligns by three bits and gives 1.65625 exactly.

\subsection{Prerequisites}\label{appendix-a1b1b4580dfd-prerequisites}

\begin{itemize}
\tightlist
\item
  \hyperref[appendix-a24fd710cf67]{Significand and hidden bit} --- Significant digit positions determine alignment and what leading digits cancel.
\item
  \hyperref[appendix-2fae058c82e1]{Exponent bias} --- Effective exponents must be decoded before comparing scale and shifting.
\item
  \hyperref[appendix-b572b631069b]{Guard round and sticky bits} --- Aligned-away information must survive sufficiently to choose the correct rounded neighbor.
\end{itemize}

\subsection{Used by}\label{appendix-a1b1b4580dfd-used-by}

\begin{itemize}
\tightlist
\item
  \hyperref[appendix-c2ac4a84c283]{Floating-point reassociation} --- uses this mechanism to establish its numerical contract.
\end{itemize}

\subsection{Recall and next depth}\label{appendix-a1b1b4580dfd-recall-and-next-depth}

Reconstruct the definition's example and name the assumptions that make it valid.

\textbf{Pending:} A full software adder must handle classes, signs, range and every rounding mode; current trace is not that implementation.

\subsection{Source and navigation}\label{appendix-a1b1b4580dfd-source-and-navigation}

\hyperref[01-foundations-i-ch001-when-arithmetic-meets-uncertainty]{M005 teaching chapter} · \hyperref[appendix-fd090d982a15-m005---parts-77-88]{Part definitions} · \href{../02-Foundations-II/main.pdf\#nameddest=appendix-4ad47bdd269a}{Dependency Map (02-Foundations-II, appendix-4ad47bdd269a)}
